\paragraph{AQ-136-DERHAM-CELLS: coherent high-weight zero-form cells still cannot give an integral full de Rham source.}
Let $k$ be a commutative ring and put $S=\Gamma_k(T)$,
with $\operatorname{wt}(T)=1$.  Suppose $P_*$ is a
simplicial graded PD $S$-algebra such that in each
simplicial degree it has a homogeneous presentation
\[
                  P_n=S\langle X_n\rangle_{\rm PD},
           \qquad \operatorname{wt}(x)\geq2
                      \quad(x\in X_n).                    \tag{136.41}
\]
The sets of variables may be infinite and may differ
with $n$.  Simplicial operators need not preserve a
chosen generating set: they are arbitrary homogeneous
PD $S$-algebra morphisms.  In particular they can encode
coherent higher relations among the cells.

Write $\Omega^1_{P/S,\mathrm{sp}}$ for the module of
\emph{special} relative PD differentials: it represents
derivations satisfying
\[
                   d\gamma_m(a)=\gamma_{m-1}(a)\,da.
                                                               \tag{136.42}
\]
Thus, in (136.41), it is the free $P_n$-module with
basis $dx$ for $x\in X_n$, of the same weights as
the variables.  This is precisely the one-form row
of the full relative PD polynomial de Rham algebra.
It is not an identification with unrestricted Beck
derivations in the algebraic theory of all PD operations.

For any prime $p$ with $k/pk\ne0$, the following
implication holds:
\[
 [T^{[p]}]=0\ \hbox{in }H_0(N(P_*^p))
 \quad\Longrightarrow\quad
 H_1\bigl(N(\Omega^1_{P_*/S,\mathrm{sp}}{}^{\!p})\bigr)
                                  \ne0.                    \tag{136.43}
\]
Here the superscript $p$ denotes weight, and $N$
denotes simplicial normalization, not a weight bound.
Consequently, if infinitely many rational primes
are nonunits in $k$, no full relative PD de Rham
theory built from (136.41) can satisfy original
condition (136.4).  In particular this rules out
that entire class of constructions over $\mathbb Z$.

\begin{proof}
Fix $p$ with $\bar k=k/pk\ne0$ and reduce the
simplicial algebra modulo $p$.  All its elements
have nonnegative weight, and its weight-zero part
is $\bar k$.  Let $\bar P_+$ denote the ideal of
positive-weight elements and define the weight-$p$
ordinary indecomposables
\[
                  Q_*=(\bar P_+/(\bar P_+)^2)^p.
                                                               \tag{136.44}
\]
This is a simplicial $\bar k$-module.  We emphasize
\emph{ordinary} indecomposables: divided powers of
positive elements are not added to the denominator.
Simplicial maps preserve the ordinary product ideal,
so no invariance of the displayed free variables
under faces is needed.

The divided-power monomial basis gives, in every
simplicial degree, the following basis of $Q_n$:
\[
       z=T^{[p]},\qquad x\in X_n\text{ with }
                                      \operatorname{wt}(x)=p.
                                                               \tag{136.45}
\]
To check this assertion completely, a monomial
involving two distinct factors among the coefficient
variable $T$ and the free variables is a product of
two positive-weight monomials, with coefficient one,
and hence is zero in (136.44).  A monomial involving
only one free variable has the form $x^{[m]}$ with
$m\operatorname{wt}(x)=p$.  Since $p$ is prime and
$\operatorname{wt}(x)\geq2$, this forces $m=1$
and $\operatorname{wt}(x)=p$.  The only remaining
coefficient monomial is $T^{[p]}$.  Any product
of two positive pure coefficient monomials of
total weight $p$ is
\[
           T^{[a]}T^{[p-a]}=\binom pa T^{[p]}=0
                    \quad(0<a<p)                         \tag{136.46}
\]
in characteristic $p$.  Products involving a free
variable cannot contribute to this coefficient
basis vector or to a single free generator.
Thus none of the vectors in (136.45) is killed or
related by the product ideal.  This proves the
asserted free $\bar k$-basis, without assuming that
$\bar k$ is a field or reduced.

There is a canonical simplicial augmentation
$\bar P_*\to\bar k$ sending every positive-weight
element to zero.  Form the simplicial module
\[
 F_*=
 \left(\bar k\otimes_{\bar P_*}
              \Omega^1_{\bar P_*/\bar S,\mathrm{sp}}\right)^p.
                                                               \tag{136.47}
\]
Its degree-$n$ term is free on the symbols $dx$
with $x\in X_n$ of weight $p$.  Composing the
relative derivation with the augmentation kills
products of positive-weight elements by the
Leibniz rule.  It therefore induces a simplicial
linear map $Q_*\to F_*$, taking $x$ to $dx$ and
$z$ to zero, because coefficients in $S$ have
zero relative differential.  The bases just
computed give an exact sequence of simplicial
$\bar k$-modules
\[
 0\longrightarrow\underline{\bar k z}
       \longrightarrow Q_*
       \longrightarrow F_*\longrightarrow0.              \tag{136.48}
\]
The first module is constant: all simplicial
operators fix the original coefficient $T^{[p]}$.

The premise of (136.43) supplies a normalized
one-chain $u\in N_1(P_*^p)$ such that
\[
                            \partial u=T^{[p]}.
                                                               \tag{136.49}
\]
Apply the relative PD derivation to $u$.  It
commutes with all simplicial maps, so
\[
 du\in N_1(\Omega^1_{P_*/S,\mathrm{sp}}{}^{\!p}),
                 \qquad\partial(du)=dT^{[p]}=0.           \tag{136.50}
\]
Reduce $u$ modulo $p$ and project it to $Q_*$;
write $v$ for its normalized image.  Then
$\partial v=z$.  Its image $a$ in $N_1(F_*)$
is a cycle.  The connecting map for (136.48)
sends its homology class to
\[
                         [z]\in H_0(\underline{\bar k z})
                                      =\bar k z,
                                                               \tag{136.51}
\]
which is nonzero.  Hence $[a]\ne0$.
Equivalently, this last step can be checked
without the long exact sequence.  If $a=\partial b$,
exactness of normalization permits lifting
$b\in N_2(F_*)$ to $\widetilde b\in N_2(Q_*)$.
Then $v-\partial\widetilde b$ belongs to
$N_1(\underline{\bar k z})=0$, whereas its
boundary is $z\ne0$, a contradiction.

The class $a$ is exactly the image of $du$ under
the chain map obtained by reduction modulo $p$
and augmentation of the differential module.
If $du$ were a boundary integrally over $k$,
its image would be a boundary in $F_*$.  The
nonzero class in (136.51) therefore proves
that (136.50) is itself a nonzero $H_1$ class.
No flat-base-change assertion for homology or
identification with an abstract cotangent complex
has been used.

Finally, condition (136.4) in cochain degree zero
would make $P_*^w$ acyclic for every sufficiently
large weight $w$, so the premise of (136.43)
would hold at every sufficiently large prime.
If infinitely many primes are nonunits in $k$,
there are arbitrarily large such primes, and
(136.43) contradicts eventual horizontal acyclicity
of the one-form row.  This proves the consequence.
\end{proof}

For a ring, $k/pk\ne0$ is equivalent to $p$ not
being a unit.  Thus this construction obstruction
is consistent with the established $p$-local and
positive-characteristic results: only finitely
many rational primes can remain nonunits in
those cases.  It is also consistent with the
standard integral theory (136.5), whose simplex
presentations have new variables of weight one.

In particular, a relative free simplicial PD
resolution of the truncated coefficient algebra
$S/(T^{[w]}:w\geq N)$, $N\geq2$, that is built
using only new variables of weights at least
$N$, cannot become admissible merely by taking
its full relative special-PD de Rham algebra.
Higher simplicial cells may repair every
zero-form relation; (136.43) then forces the
one-form obstruction.  This rules out that
coherent repair, not all possible coherent
repairs: weight-one variables, additional form
relations, or other non-free theories are
outside the theorem.  The original integral
common-source and localization questions remain
separate.
